\section{Application: Wind Farm Layout Optimization}
\label{sec:application}

\paragraph{Wake model.}
We use the Bastankhah Gaussian wake deficit
model~\cite{bastankhah2014new}, implemented in \texttt{pixwake}, a
JAX-native wake model library. The wake superposition uses a
fixed-point iteration with a custom VJP rule, enabling
\texttt{jax.grad} to compute exact gradients of AEP with respect to
turbine positions. The wake expansion coefficient is fixed at
$k=0.04$ for all experiments. The LLM receives the simulation
object and computes gradients freely but cannot modify wake
parameters. \texttt{pixwake} has not been cross-validated against
FLORIS or PyWake for this benchmark; systematic wake-model biases
could affect absolute AEP values but not relative comparisons
between methods.

\paragraph{Training farm: DEI~1.}
50 IEA 15\,MW turbines~\cite{gaertner2020iea15mw}
($D=240$\,m), 960\,m minimum spacing ($4D$), 5-vertex convex
polygon ($\sim$$20\!\times\!35$\,km, generously sized relative to
the turbine count), 24-sector wind rose from 10 years of observed
DEI data~\cite{energinet2022dei}.

\paragraph{Held-out farm: IEA Wind ROWP.}
74 IEA 10\,MW turbines~\cite{borrassunyeriea2024rowp,bortolotti2019iea10mw}
($D=198$\,m), 792\,m minimum spacing, 4-vertex elongated polygon
($\sim$$19\!\times\!16$\,km), 12-sector Weibull wind rose. ROWP
differs from DEI on every axis an optimizer might exploit:
turbine count $+48\%$, rotor $-17.5\%$, spacing $-17.5\%$, polygon
geometry, wind rose. These differences reflect real variation in
offshore planning cases, not artificial hardening.

\paragraph{Generalization matrix.}
To test whether discovered schedules generalize beyond the
training configuration, we construct a $2 \times 6 \times 4 = 48$-cell
evaluation matrix (Table~\ref{tab:matrix}). Each cell is a distinct
layout problem with its own 500-multi-start baseline. The two farm
polygons (DEI and ROWP) provide a single held-out polygon for
transfer testing; six turbine
counts ($N \in \{30, \ldots, 80\}$) test scaling under increasing
packing density; and four wind roses isolate directional effects:
(i)~\emph{uniform}---a single wind direction retaining DEI's speed
distribution; (ii)~\emph{omnidirectional}---24 equally-weighted
directions; (iii)~\emph{DEI rose}---the training distribution;
(iv)~\emph{ROWP rose}---the held-out distribution. The wind speed
marginal is preserved across all roses so that only directional
structure varies.

\paragraph{Constraint handling in the baseline solver.}
\texttt{topfarm\_sgd\_solve} uses differentiable boundary and
spacing penalties added to the objective with an adaptive
multiplier $\alpha_t = \alpha_0 \cdot (\eta_0 / \eta_t)$ that
grows as the learning rate decays. This 1/$\eta$-coupling allows
large infeasible moves early and enforces feasibility as the
optimizer settles---and is the mechanism that no LLM-generated
custom optimizer reproduces in full-optimizer mode (Sec.~\ref{sec:results}).

\paragraph{Models, time budget, hot-start.}
Each run has a 5-hour wall-clock budget; per-evaluation timeout
60\,s (schedule-only) or 180\,s (full-optimizer). The schedule-only
timeout is a \emph{compute equalizer}: every schedule runs the same
8000-step Adam loop, so all evaluations take identical wall-clock
time ($\sim$30\,s at $N\!=\!50$). In full-optimizer mode, different
approaches have different compute profiles---a single SLSQP run
with JAX JIT compilation takes $\sim$40\,s, leaving room for at
most 2--3 restarts at 180\,s.
We evaluate two frontier LLM agents through their rich CLI
scaffolds: Claude Code~\cite{anthropic2025claude} and Gemini~2.5
Flash via Gemini CLI~\cite{gemini2025flash}. Both runs hot-start
from a single-start \texttt{topfarm\_sgd\_solve} seed scoring below
the 500-start baseline. Wake simulations are deterministic (JAX
\texttt{float64}); LLM sampling is stochastic.
